\section*{Capítulo \ref{ch:g12:synthesis-strategy} --- Estratégia de síntese e química verde}

\begin{solution}{exo:g12:synthesis-strategy:1}
Hidratação: $46.0/(28.0 + 18.0) = \qty{100}{\percent}$. Fermentação:
$2 \times 46.0/180.0 = \qty{51.1}{\percent}$.
\end{solution}

\begin{solution}{exo:g12:synthesis-strategy:2}
Os dois grupos, o \ce{Z} na \omterm{def:g11:functional-groups:amine}{amina} da glicina e o \omterm{def:g11:functional-groups:carboxylic-acid}{éster} metílico no ácido
da alanina, são colocados na etapa 1 e retirados na etapa 3.
\end{solution}

\begin{solution}{exo:g12:synthesis-strategy:3}
Princípio 5: usar \omterm{def:g6:solutions-and-solubility:solute}{solventes} e condições de reação mais seguros.
\end{solution}

\begin{solution}{exo:g12:synthesis-strategy:4}
Um excesso de etanol torna $Q < K$, logo a reação avança mais no
sentido direto. Um \omterm{def:g12:catalysis:catalyst}{catalisador} acelera os dois sentidos da mesma forma:
chega-se ao mesmo \omterm{def:g10:reaction-progress-table:system}{estado final}, só que mais cedo.
\end{solution}

\begin{solution}{exo:g12:synthesis-strategy:5}
Sim: só a \omterm{def:g11:functional-groups:carbonyl}{cetona} é reduzida, a um grupo \omterm{def:g11:functional-groups:alcohol}{álcool} secundário; o \omterm{def:g11:functional-groups:carboxylic-acid}{éster} não
muda.
\end{solution}

\begin{solution}{exo:g12:synthesis-strategy:6}
Mais \omterm{def:g11:functional-groups:carboxylic-acid}{éster}: um excesso de um \omterm{def:g7:chemical-reactions:reactant}{reagente}, ou retirar a água (separador de
água) à medida que ela se forma. Mais depressa: aquecer sob refluxo,
acrescentar um \omterm{def:g12:catalysis:catalyst}{catalisador} ácido.
\end{solution}

\begin{solution}{exo:g12:synthesis-strategy:7}
Com \ce{NaBH4}:
\[ \ce{CH3-CH(OH)-CH2-CH2-CO-O-CH3} . \]
Com \ce{LiAlH4}:
\[ \ce{CH3-CH(OH)-CH2-CH2-CH2OH} \]
(pentano-1,4-diol), mais metanol vindo da outra metade do \omterm{def:g11:functional-groups:carboxylic-acid}{éster}.
\end{solution}

\begin{solution}{exo:g12:synthesis-strategy:8}
A adição (\qty{100}{\percent}), a esterificação
(\qty{83}{\percent}) e a rota de três etapas do ibuprofeno
(\qty{77}{\percent}). $77.4/40.0 = 1.9$: quase duas vezes melhor.
\end{solution}

\begin{solution}{exo:g12:synthesis-strategy:9}
$180.0/(138.0 + 102.0) = \qty{75.0}{\percent}$. Ácido etanoico:
$60.0/180.0 = \qty{0.333}{kg}$ por quilograma de aspirina.
\end{solution}

\begin{solution}{exo:g12:synthesis-strategy:10}
O aquecimento acelera a reação; o condensador devolve os vapores ao
balão, de modo que nenhum \omterm{def:g7:chemical-reactions:reactant}{reagente} ou \omterm{def:g6:solutions-and-solubility:solute}{solvente} se perde. O aquecimento
melhora a velocidade, não o \omterm{def:g10:synthesis-yield:yield}{rendimento} final.
\end{solution}

\begin{solution}{exo:g12:synthesis-strategy:11}
Seis etapas contra três. Princípios: prevenir os resíduos (1), maximizar
a \omterm{def:g12:synthesis-strategy:atom-economy}{economia atômica} (2), usar \omterm{def:g12:catalysis:catalyst}{catalisadores} (9); e também menos etapas,
menos energia (6).
\end{solution}

\begin{solution}{exo:g12:synthesis-strategy:12}
$0.80 \times 0.70 \times 0.90 = 0.504$: \qty{50.4}{\percent}.
\omterm{def:g5:raw-materials-and-recycling:raw-material}{Matéria-prima}: $0.10/0.504 = \qty{0.20}{mol}$.
\end{solution}

\begin{solution}{exo:g12:synthesis-strategy:13}
Proteger a \omterm{def:g11:functional-groups:amine}{amina} da alanina (\ce{Z}) e o ácido da glicina (\omterm{def:g11:functional-groups:carboxylic-acid}{éster}
metílico); formar a ligação \omterm{def:g11:functional-groups:amine}{amida} entre o ácido livre da alanina e a
\omterm{def:g11:functional-groups:amine}{amina} livre da glicina; retirar os dois \omterm{def:g12:synthesis-strategy:protecting-group}{grupos protetores}. Três fases,
isto é, cinco reações se cada proteção e cada desproteção for contada.
\end{solution}

\begin{solution}{exo:g12:synthesis-strategy:14}
O \omterm{def:g12:catalysis:catalyst}{catalisador} aplica o princípio 9. Mas os \qty{250}{\celsius} vão
contra o princípio 6, o benzeno contra os princípios 3 e 5, e uma
\omterm{def:g12:synthesis-strategy:atom-economy}{economia atômica} de \qty{35}{\percent} contra o princípio 2:
$65/35 = \qty{1.9}{kg}$ de resíduos por quilograma de produto. Uma
característica verde não faz um processo verde.
\end{solution}

\begin{solution}{exo:g12:synthesis-strategy:15}
% ledger: crit:CO2-T, crit:CO2-p
\qty{40}{\celsius} e \qty{100}{bar} estão acima do ponto crítico do
dióxido de carbono, \qty{31}{\celsius} e \qty{73.8}{bar}. A pressão é
cem vezes a pressão atmosférica: o vaso é de aço grosso. O dióxido de
carbono não é tóxico nem inflamável, volta a ser gás sem deixar resíduo,
e pode ser reutilizado; nenhum \omterm{def:g6:solutions-and-solubility:solute}{solvente} clorado escapa.
\end{solution}

\begin{solution}{pb:g12:synthesis-strategy:1}
% ledger: ibu:bhc-ae-recovered
\textbf{1.} C: $10 + 4 + 1 = 15 = 13 + 2$; H: $14 + 6 + 2 = 22 = 18 + 4$;
O: $3 + 1 = 4 = 2 + 2$.

\textbf{2.} O ácido etanoico, \ce{C2H4O2}.

\textbf{3.} O nitrogênio, o cloro e o sódio: o ibuprofeno só contém
carbono, hidrogênio e oxigênio.

\textbf{4.} A mais nova: princípio 9, \omterm{def:g12:catalysis:catalyst}{catalisadores} em vez de \omterm{def:g7:chemical-reactions:reactant}{reagentes}
estequiométricos.

\textbf{5.} $13 \times 12.0 + 18 \times 1.0 + 2 \times 16.0 =
\qty{206.0}{g/mol}$.

\textbf{6.} $134.0 + 102.0 + 122.5 + 68.0 + 19.0 + 33.0 + 36.0 =
\qty{514.5}{g/mol}$.

\textbf{7.} $206.0/514.5 = \qty{40.0}{\percent}$.

\textbf{8.} $134.0 + 102.0 + 2.0 + 28.0 = \qty{266.0}{g/mol}$.

\textbf{9.} $206.0/266.0 = \qty{77.4}{\percent}$.

\textbf{10.} $(206.0 + 60.0)/266.0 = \qty{100}{\percent}$ no papel; mais
de \qty{99}{\percent} na prática.

\textbf{11.} $514.5/206.0 = \qty{2.50}{t}$.

\textbf{12.} $2.50 - 1 = \qty{1.50}{t}$.

\textbf{13.} $266.0/206.0 = \qty{1.29}{t}$.

\textbf{14.} $60.0/206.0 = \qty{0.29}{t}$.

\textbf{15.} $1.50 - 0.29 = \qty{1.21}{t}$.

\textbf{16.} $0.90^6 = 0.53$: \qty{53}{\percent}.

\textbf{17.} $0.90^3 = 0.73$: \qty{73}{\percent}.

\textbf{18.} Cada etapa usa \omterm{def:g6:solutions-and-solubility:solute}{solventes}, separações e purificações, e perde
um pouco de produto: menos etapas, menos dessas perdas.

\textbf{19.} Quase nenhum: o seu único subproduto é usado.

\textbf{20.} A rota de três etapas evita cerca de \qty{1.5}{t} de
resíduos por tonelada de ibuprofeno (\qty{1.2}{t} mesmo que o seu ácido
etanoico fosse jogado fora).
\end{solution}
